\documentclass[10pt,a4paper]{amsart}
\usepackage[margin=19mm]{geometry}
\usepackage{amsmath,amssymb,mathtools}
\usepackage[scr=rsfs,cal=euler]{mathalfa}
\usepackage{xfrac}
\usepackage[unicode,hidelinks,bookmarks=false]{hyperref}
\newcommand{\ie}{i.e.\ }
\newcommand{\cf}{cf.\ }
\newcommand{\resp}{respectively\ }
\newcommand{\Subsection}{\subsection}
\providecommand{\nobreakdash}{\nobreak\hbox{-}}
\setlength{\emergencystretch}{2em}
\multlinegap0pt
\usepackage{bm}
\usepackage{bbm}
\usepackage{dsfont}
\usepackage{xspace}
\usepackage{cancel}
\usepackage{mathtools}
\usepackage{amsthm}
\makeatletter
\@ifundefined{thm}{\newtheorem{thm}{Thm}}{}
\@ifundefined{claim}{\newtheorem{claim}[thm]{Claim}}{}
\@ifundefined{lem}{\newtheorem{lem}[thm]{Lemma}}{}
\@ifundefined{prop}{\newtheorem{prop}[thm]{Proposition}}{}
\makeatother
\newcommand{\score}{\mathrm{score}}
\newcommand{\sig}{\sigma}
\newcommand{\sm}{\setminus}
\newcommand{\sub}{\subseteq}
\newcommand{\tour}{\vec{S}}
\renewcommand{\c}[1]{\mathcal{#1}}
\title[Metric Dimensions of March Madness Brackets]{Metric Dimensions of March Madness Brackets}
\author{Sam Spiro}
\date{Source: arXiv:2602.06950v1 (CC BY 4.0)}
\begin{document}
\maketitle

\noindent\emph{Source and licence.} Metric Dimensions of March Madness Brackets. Version arXiv:2602.06950v1 (CC BY 4.0), distributed under the Creative Commons Attribution 4.0 International licence, as recorded in the arXiv licence metadata for this exact version and shown on the abstract page https://arxiv.org/abs/2602.06950v1. Full rights notice: licenses/arxiv-2602.06950-CC-BY-4.0.txt.\par\smallskip

\noindent\textit{Editorial scope.} Counted unit: the Lem whose source label is \texttt{favorite team} in the article (source0.tex lines 651--653, source label \texttt{favorite team}), delivered together with the article's own complete original proof of it (source0.tex lines 654--696, one \texttt{proof} environment of 43 lines). No mathematics is written, repaired or re-derived: statement and proof are the article's own text line for line. Required context retained uncounted: Pu Sets (lines 328--337); bracket facts (lines 361--367); a implies b (lines 642--644); Claim A size (lines 749--751). Every definition, display and label of those lines is retained unchanged. Omitted from this package and not redistributed: 1--327, 338--360, 368--641, 645--650, 697--748, 752--992. Nothing inside a retained excerpt is omitted or altered: every retained body is delivered exactly as the article's lines are written, so no omission receipt of any kind accompanies this package. Citations: the retained excerpts carry no citation command at all, so no bibliography entry of the article is transcribed and no reference is redistributed. Reference adaptation: none was needed and none was made; every reference inside the delivered unit resolves inside it, so no unresolved reference or citation is produced. Printed slips kept literal: no printed word, symbol, hypothesis or number of the counted statement or of its proof was altered. Layout and class: the article's own class is \texttt{article} with packages including amssymb, amsfonts, amscd, amsthm, xy, enumerate, bbm, mathrsfs, tikz, geometry, mathtools, hyperref, cleveref, color; the wrapper is the lane's pinned \texttt{amsart} editorial wrapper at 10pt on A4, which restates the theorem-like environments the retained text needs and adds the article's own macro definitions that the retained text needs (\texttt{\textbackslash c}, \texttt{\textbackslash score}, \texttt{\textbackslash sig}, \texttt{\textbackslash sm}, \texttt{\textbackslash sub}, \texttt{\textbackslash tour}), with extra wrapper packages (bm, bbm, dsfont, xspace, cancel, mathtools). The delivered numbering is the unit's own. Assets: no retained span contains a figure environment, a graphics-inclusion command, a PDF-page inclusion, a table or a TikZ picture; no image file is copied into this package, the package declares no asset dependency and no image file is redistributed with it. Licence: version arXiv:2602.06950v1 of this article is distributed under the Creative Commons Attribution 4.0 International licence, as recorded in the arXiv licence metadata for this exact version and shown on the abstract page https://arxiv.org/abs/2602.06950v1. A full rights notice is delivered at licenses/arxiv-2602.06950-CC-BY-4.0.txt. Characters: every retained excerpt is pure ASCII, so the delivered engine, XeLaTeX with its default Latin Modern text font, reports no missing character.
\begin{prop}\label{Pu Sets}
	Let $\tour$ be a single-elimination tournament.
	\begin{itemize}
		\item[(i)] Every $u\in V(\tour)$ has $P(u)\ne \emptyset$.
		\item[(ii)] For all $u,v\in V(\tour)$ we have $P(u)= P(v)$ if and only if $u= v$
		\item[(iii)] We have $P(u)\sub P(v)$ for $u,v\in V(\tour)$ if and only if there exists a directed walk from $u$ to $v$.
		\item[(iv)] For all distinct $u,v\in V(\tour)$ we either have $P(u)\cap P(v)=\emptyset,\ P(u)\subsetneq P(v)$, or $P(v)\subsetneq P(u)$.
		\item[(v)] For every match $x\in M(\tour)$ the sets $\{P(u):u\in N^-(x)\}$ partition $P(x)$ and satisfy that $P(u)\subsetneq P(x)$ and $P(u)\ne \emptyset$ for all $u\in N^-(x)$.
	\end{itemize}
\end{prop}

\begin{lem}\label{bracket facts}
	Let $\tour$ be a single-elimination tournament and $B$ a bracket.
	\begin{itemize}
		\item[(i)] We have $B(u)\in P(u)$ for all $u \in V(\tour)$.
		\item[(ii)] For each $v\in V(\tour)$, we have $B(u)=B(v)$ for every $u\in V(\tour)$ with $P(u)\sub P(v)$ and $B(v)\in P(u)$.
	\end{itemize}
\end{lem}

\begin{lem}\label{a implies b}
	Let $\tour$ be a single-elimination tournament on at least 2 vertices, $a\in P(\tour)$, and $x_a$ the unique out-neighbor of $a$.  If $b\in P(x_a)\sm \{a\}$, then $b\in P(x)$ for every match $x$ with $a\in P(x)$.
\end{lem}

	\begin{claim}\label{Claim A size}
		The number of 2-element subsets of $\c{A}_2$ is $\sum_{x\in M(\tour)\sm V(\tour_u)} \lfloor \frac{|N^-(x)\cap P(\tour)\sm P(u)|}{2} \rfloor$.
	\end{claim}

\begin{lem}\label{favorite team}
	Let $\tour$ be a single-elimination tournament with at least 2 vertices and $\widehat{B}$ a bracket.  If $a,b$ are distinct players such that the out-neighbor $x_a\in N^+(a)$ has $b\in \{\widehat{B}(v):v\in N^-(x_a)\}$, then for any brackets $B,B'$ and scoring system $\sig$ satisfying that $\score_\sig(\widehat{B}_a,B)=\score_\sig(\widehat{B}_a,B')$ and $\score_\sig(\widehat{B}_b,B)=\score_\sig(\widehat{B}_b,B')$, we have for all $x\in M(\tour)$ that $B(x)=a$ if and only if $B'(x)=a$ .
\end{lem}

\begin{proof}
	We begin with a few observations about our setup.
	\begin{claim}\label{Claim favorite team nice vertex}
		There exists an in-neighbor $v\in N^-(x_a)$ with $\widehat{B}(v)=b$ and with $b\in P(v)$ and $a\notin P(v)$.
	\end{claim}
	\begin{proof}
		The existence of $v\in N^-(x_a)$ with $\widehat{B}(v)=b$ is by hypothesis, and this immediately implies $b\in P(v)$ by Lemma~\ref{bracket facts}(i).  Moreover, because $a\in N^-(x_a)$ (with necessarily $a\ne v$ since $b\notin P(a)$), we must have $a\notin P(v)$ by Proposition~\ref{Pu Sets}(v).
	\end{proof}
	From now on any reference to $v$ refers to the vertex guaranteed by Claim~\ref{Claim favorite team nice vertex}.
	
	\begin{claim}\label{Claim favorite team a implies b}
		Every match $x$ with $a\in P(x)$ has $b\in P(x)$.
	\end{claim}
	\begin{proof}
		By Claim~\ref{Claim A size} and Proposition~\ref{Pu Sets}(v) we have $b\in P(v)\sub P(x_a)\sm \{a\}$, so the claim follows from Lemma~\ref{a implies b}.
		
		
		%For any match $x$ with $a\in P(x)$ there exists a directed walk $(u_1,\ldots,u_t)$ from $a$ to $x$ by definition.  Note that $u_2=x_a$ since $x_a$ is the unique out-neighbor of $a$.  This means $(u_2,\ldots,u_t)$ is a directed walk from $x_a$ to $a$, implying $P(x_a)\sub P(x)$ by Proposition~\ref{Pu Sets}.  But because $v\in N^-(x_a)$ and $b\in P(v)\sub P(x)$ by Proposition~\ref{Pu Sets}(v) we must also have $b\in P(x)$, proving the claim.
	\end{proof}
	
	\begin{claim}
		We have $\widehat{B}_a(x)=\widehat{B}_b(x)$ for every match $x$ with $a\notin P(x)$.
	\end{claim}
	\begin{proof}
		If $x$ is a match with $a,b\notin P(x)$ then $\widehat{B}_a(x)=\widehat{B}(x)=\widehat{B}_b(x)$ by definition, so it suffices to prove this for matches with $a\notin P(x)$ and $b\in P(x)$.  Because $\widehat{B}(v)=b$, we must have $\widehat{B}(x)=b$ for all $x$ with $P(x)\sub P(v)$ by Lemma~\ref{bracket facts}(ii), which means $\widehat{B}_a(x)=\widehat{B}(x)=b=\widehat{B}_b(x)$ for all such $x$.  The only remaining case to check then is when $P(x)\not \sub P(v)$.  Because $b\in P(x)\cap P(v)$, any such $x$ has $P(v)\subsetneq P(x)$ by Proposition~\ref{Pu Sets}(iv), so by Proposition~\ref{Pu Sets}(iii) there is a directed walk $(u_1,\ldots,u_t)$ from $v$ to $x$.  Note that $u_2=x_a$ since this is the unique out-neighbor of $v$, so $(u_2,\ldots,u_t)$ is a directed walk from $x_a$ to $x$, implying by Proposition~\ref{Pu Sets}(iii) that $a\in P(x_a)\sub P(x)$, a contradiction to the assumption $a\notin P(x)$.  We conclude the result.
	\end{proof}
	Let $M_a\sub M(\tour)$ denote the set of matches $x$ with $a\in P(x)$.  From this claim and the definition of $\score_\sig$, we see that for any bracket $B^*$ we have
		\begin{equation}\score_\sig(\widehat{B}_a,B^*)-\score_\sig(\widehat{B}_b,B^*)=\sum_{x\in M_a:B^*(x)=a}\sig(x)-\sum_{x\in M_a:B^*(x)=b}\sig(x).\label{eq:difference in scores}\end{equation}

	Now assume for contradiction that the lemma is false for some match $y$ and choose this $y$ with $|P(y)|$ as small as possible.  Without loss of generality we may assume that $B(y)=a$ and $B'(y)\ne a$.
	
	We claim that $B(x)\ne b$ for every $x\in M_a$.  Indeed if $B(x)=b$ for some $x\in M_a$ then this means $a,b\in P(x)$ by Lemma~\ref{bracket facts}(i) and the definition of $M_a$.  Similarly we have $a,b\in P(y)$ due to Lemma~\ref{bracket facts}(i) and Claim~\ref{Claim favorite team a implies b}.  By Proposition~\ref{Pu Sets}(iv) we have that either $P(x)\sub P(y)$ or $P(y)\sub P(x)$.  In the first case because $a\in P(x)\sub P(y)$ and $B(y)=a$ we must have $B(x)=a$ by Lemma~\ref{bracket facts}(ii), a contradiction.  The same argument holds if $P(y)\sub P(x)$, proving the claim.  With this, we have by \eqref{eq:difference in scores} that
	\[\score_\sig(\widehat{B}_a,B)-\score_\sig(\widehat{B}_b,B)= \sum_{x\in M_a:B(x)=a}\sig(x)\ge \sum_{x\in M_a: P(x)\sub P(y)} \sig(x),\]
	with this last step using Lemma~\ref{bracket facts}(ii) to ensure $B(x)=a$ for all $x\in M_a$ with $P(x)\sub P(y)$.
	
	We claim that $B'(x)\ne a$ for any $x\in M_a$ with $P(x)=P(y)$ or $P(x)\not\sub P(y)$.  Indeed, the first part follows since Proposition~\ref{Pu Sets}(ii) implies $P(x)=P(y)$ only can happen for $x=y$ and we assumed $B'(y)\ne a$.  For the second part, having $x\in M_a$ means $a\in P(x)\cap P(y)$, so $P(x)\not \sub P(y)$ would imply $P(y)\sub P(x)$ by Proposition~\ref{Pu Sets}(iv).  But now $B'(x)=a$ would imply $B'(y)=a$ by Lemma~\ref{bracket facts}(ii), a contradiction.  This proves the claim, and hence by \eqref{eq:difference in scores} we have 
	\[\score_\sig(\widehat{B}_a,B')-\score_\sig(\widehat{B}_b,B')\le \sum_{x\in M_a:B'(x)=a}\sig(x)= \sum_{x\in M_a: P(x)\subsetneq P(y)}\sig(x)< \sum_{x\in M_a: P(x)\sub P(y)} \sig(x),\]
	with this last step using that $\sig(y)>0$ by definition of scoring systems.
	
	In total we find that 
	\[\score_\sig(\widehat{B}_a,B)-\score_\sig(\widehat{B}_b,B)>\score_\sig(\widehat{B}_a,B')-\score_\sig(\widehat{B}_b,B'),\]
	which implies that either  $\score_\sig(\widehat{B}_a,B)\ne \score_\sig(\widehat{B}_a,B')$ or $\score_\sig(\widehat{B}_b,B)\ne \score_\sig(\widehat{B}_b,B')$, a contradiction to our hypothesis.  We conclude the result.
\end{proof}

\end{document}
